\clearpage\fancyhead[L]{\small VESTA.}\hypertarget{native-366-lower}{}\label{native-366-lower}
\begin{center}\textbf{VESTA.}\par\vspace{4pt}\hrule\vspace{4pt}\fontsize{10}{12}\selectfont Göttingische gelehrte Anzeigen. No. 205. pp. 2041\;..\;2045. 24 December 1812.\par\vspace{4pt}\hrule\vspace{4pt}\end{center}

\par\indent In October of this year, Vesta, the most recently discovered of the new planets, came into opposition with the Sun for the fourth time since its discovery. The fourth opposition is always especially important for the theory of a new planet, since it first makes it possible to determine the orbit from oppositions alone. So Vesta was observed on the meridian at our observatory around the time of opposition as often as the poor weather allowed. Professor \textsc{Gauss} paid particular attention to the declinations. The right \SourceLeaf{367}ascensions were therefore measured only in passing, since there was reason to hope that other observatories could provide better measurements than the mural quadrant. It was all the more necessary to limit the observations to what the mural quadrant could measure well because the planet was considerably less than half as bright as at the previous year’s opposition. It was therefore often difficult to observe this time with the instrument, which gathered little light. Vesta was observed five times in all. The observations on 24 and 27 October, however, were closer to estimates because atmospheric conditions were very poor. The three declinations on 25 October, 28 October and 1 November should be regarded as very good. On 1 November the right ascension was measured only incidentally, at a single wire.

\medskip
\Needspace{12\baselineskip}\noindent\emph{Observation of Vesta in Göttingen.}

\begin{center}
\fontsize{10}{12}\selectfont\setlength{\tabcolsep}{4pt}\renewcommand{\arraystretch}{1.15}\begin{tabular}{|c|c|c|c|}
\hline
1812 & Mean time & Right ascension & Northern declination \\
\hline
Oct. 24 & $12^{\mathrm u}\;3^{\mathrm m}\;39^{\mathrm s}$ & . . . . & $1^\circ\,54'\,49''$ \\
25 & $11\;58\;43$ & $33^\circ\,52'\,33''1$ & $1\;50\;15.8$ \\
27 & $11\;48\;53$ & . . . . & $1\;41\;56$ \\
28 & $11\;43\;58$ & $33\;\;8\;\;1.8$ & $1\;38\;40.9$ \\
Nov. 1 & $11\;24\;21$ & $32\;\;9\;35.9$ & $1\;25\;\;0.9$ \\
\hline
\end{tabular}
\end{center}

\Needspace{13\baselineskip}\par\indent At the Seeberg Observatory the planet could be observed only twice. Mr. von \textsc{Lindenau} communicated the following observations to Professor \textsc{Gauss}:

\begin{center}
\fontsize{10}{12}\selectfont\setlength{\tabcolsep}{4pt}\renewcommand{\arraystretch}{1.15}\begin{tabular}{|c|c|c|c|}
\hline
1812 & Mean time & Right ascension & Northern declination \\
\hline
Oct. 25 & $11^{\mathrm u}\;58^{\mathrm m}\;46^{\mathrm s}8$ & $33^\circ\,52'\,31''7$ & $1^\circ\,50'\,28''1$ \\
Nov. 1 & $11\;24\;23.1$ & $32\;\;9\;11.0$ & $1\;25\;28{,}9$ \\
\hline
\end{tabular}
\end{center}

\Needspace{17\baselineskip}\par\indent Finally, Mr. von \textsc{Zach} made the following meridian observations at his observatory near Marseille:

\begin{center}
\fontsize{10}{12}\selectfont\setlength{\tabcolsep}{4pt}\renewcommand{\arraystretch}{1.15}\begin{tabular}{|c|c|c|c|}
\hline
1812 & Mean time & Right ascension & Northern declination \\
\hline
Oct. 21 & $12^{\mathrm u}\;18^{\mathrm m}\;21^{\mathrm s}1$ & $34^\circ\,51'\,13''2$ & $2^\circ\,\;7'\,57''5$ \\
23 & $12\;\;8\;32.3$ & $34\;21\;53.4$ & $1\;54\;39{,}9$ \\
24 & $12\;\;3\;37.5$ & $34\;\;7\;\;8.4$ & $1\;54\;25.3$ \\
26 & $11\;53\;47.5$ & $33\;37\;31.2$ & $1\;46\;15.7$ \\
28 & $11\;43\;57.4$ & $33\;\;7\;52.2$ & $1\;38\;54.5$ \\
29 & $11\;39\;\;2.4$ & $32\;53\;\;2.4$ & $1\;35\;\;0.1$ \\
\hline
\end{tabular}
\end{center}

\par\indent All these observations were compared with the latest elements, communicated in the Monatl. Corresp., vol. XXIV, p. 502. To obtain better agreement, the epoch was increased by $12'\,47''3$, giving the following differences.

\medskip
\Needspace{14\baselineskip}\noindent\emph{Observations in Göttingen.} 

\begin{center}
Difference

\fontsize{10}{12}\selectfont\setlength{\tabcolsep}{7pt}\renewcommand{\arraystretch}{1.15}\begin{tabular}{|c|c|c|}
\hline
 & Right ascension & Declination \\
\hline
Oct. 24 & . . . & $-38''68$: \\
25 & $+3''04$ & $-13.26$ \\
27 & . . . & $+16.10$: \\
28 & $+6.37$ & $-13.32$ \\
Nov. 1 & $-15.16$ & $-10.17$ \\
\hline
\end{tabular}
\end{center}

\Needspace{12\baselineskip}\SourceLeaf{368}\noindent\emph{Observations at the Seeberg Observatory.}

\begin{center}
Difference.

\fontsize{10}{12}\selectfont\setlength{\tabcolsep}{7pt}\renewcommand{\arraystretch}{1.15}\begin{tabular}{|c|c|c|}
\hline
 & Right ascension & Declination \\
\hline
Oct. 25 & $+6''44$ & $-23''41$ \\
Nov. 1 & $+11.71$ & $-37.70$ \\
\hline
\end{tabular}
\end{center}

\medskip
\Needspace{15\baselineskip}\noindent\emph{Observations at Lacapellete near Marseille.}

\begin{center}
Difference.

\fontsize{10}{12}\selectfont\setlength{\tabcolsep}{7pt}\renewcommand{\arraystretch}{1.15}\begin{tabular}{|c|c|c|}
\hline
 & Right ascension & Declination \\
\hline
Oct. 21 & $+2''73$ & $-52''74$ \\
23 & $+3.93$ & $+220.85$ \\
24 & $+3.80$ & $-17.41$ \\
26 & $+5.06$ & $-13.51$ \\
28 & $+4.80$ & $-29.19$ \\
29 & $+7.20$ & $-11.60$ \\
\hline
\end{tabular}
\end{center}

\par\indent The right ascensions can be seen to agree well, while the declinations agree less closely. Professor \textsc{Gauss} felt he should rely chiefly on the three declinations observed here, which agree very well with one another. There appears to have been an error of 4 minutes in Mr. von \textsc{Zach}’s observation on 23 October. On that assumption, at least the four observations on 23, 24, 26 and 29 October agree fairly well with one another and agree closely with those made here. The following result was finally obtained for the opposition:

\medskip
\begin{center}\fontsize{10}{12}\selectfont\setlength{\tabcolsep}{2pt}\renewcommand{\arraystretch}{1.15}\begin{tabular}{l}
1812 October $8^{\mathrm u}\;54^{\mathrm m}\;44^{\mathrm s}$\\Mean time in Göttingen
\end{tabular}\end{center}
\begin{center}\fontsize{10}{12}\selectfont\setlength{\tabcolsep}{2pt}\renewcommand{\arraystretch}{1.15}\begin{tabular}{lcr}
True longitude & $.\;.\;.\;.\;.\;.\;.\;.\;.\;.\;.$ & $32^\circ\;17'\;40''8$\\
Geocentric latitude & $.\;.\;.\;.\;.\;.\;.\;.\;.$ & $11\;5\;32.3\;\text{S.}$\\
\end{tabular}\end{center}

\Needspace{22\baselineskip}\par\indent Professor \textsc{Gauss} then entrusted Mr. \textsc{Enke} with correcting the elements on the basis of the four oppositions observed so far. Mr. \textsc{Enke} is studying astronomy here with outstanding success and is already highly skilled in astronomical calculation. His result is as follows:

\medskip
\begin{center}\fontsize{10}{12}\selectfont\setlength{\tabcolsep}{2pt}\renewcommand{\arraystretch}{1.15}\begin{tabular}{l}
Epoch of longitude, 1814 February 13.\\$0^{\mathrm u}$ in Göttingen $\; .\; 154^\circ\;48'\;29''4$
\end{tabular}\end{center}
\begin{center}\fontsize{9}{11}\selectfont\setlength{\tabcolsep}{2pt}\renewcommand{\arraystretch}{1.15}\begin{tabular}{lcr}
Daily mean tropical motion & $.\;.\;.\;.\;.\;.$ & $978''1642$\\
Longitude of perihelion & $.\;.\;.\;.\;.\;.\;.\;.\;.\;.\;.\;.$ & $249^\circ\;38'\;31.4$\\
Longitude of the ascending node & $.\;.\;.\;.\;.\;.\;.\;.$ & $103\;12\;25.3$\\
Inclination of the orbit & $.\;.\;.\;.\;.\;.\;.\;.\;.\;.\;.\;.\;.$ & $7\;8\;5.0$\\
Eccentricity & $.\;.\;.\;.\;.\;.\;.\;.\;.$ & $=0.0894779=\sin 5\;8\;0.86$\\
Logarithm of the semimajor axis & $.\;.\;.\;.\;.\;.\;.$ & $0.3731047$\\
\end{tabular}\end{center}

\Needspace{12\baselineskip}\par\indent According to Mr. \textsc{Enke}’s calculation from these elements, the next opposition will occur on 13 February 1814 at $10^{\mathrm u}\;21^{\mathrm m}\;26^{\mathrm s}$, at longitude $144^\circ\,37'\,51''$ and latitude $8^\circ\,2'\,19''$ N., when the planet’s brightness will be $= 0.08698$. At the first four oppositions it was

\medskip
\begin{center}\begin{tabular}{ll}$1808\;.\;.\;0.09516,$ & $1810\;.\;.\;0.06440,$\\$1811\;.\;.\;0.16540,$ & $1812\;.\;.\;0.06641$\end{tabular}\end{center}

\medskip
\noindent A complete ephemeris for 1813 and 1814, calculated by Mr. \textsc{Enke}, will appear elsewhere.

\Needspace{39\baselineskip}\SourceLeaf{369}\par\indent We take this opportunity to give a brief account of an important piece of work that Mr. \textsc{Nicolai} has just completed on the second comet of 1811. This comet has already been discussed twice in these pages (1811, no. 201, and 1812, no. 15 [pp. 342 and 345 of this volume]). The latter issue reported the results of Mr. \textsc{Nicolai}’s calculations. The exceptional precision that observations of this comet allowed, and the appreciable span they covered, seemed to Professor \textsc{Gauss} sufficient to reach some conclusion about the ellipticity of the orbit. He therefore encouraged Mr. \textsc{Nicolai} to undertake this interesting investigation, and the outcome fully confirmed that judgment. Of all comets whose orbital periods have been calculated from a single apparition, this one surely ranks among those whose periods are best determined. Mr. \textsc{Nicolai}’s result is as follows:

\Needspace{22\baselineskip}\medskip
\begin{center}\fontsize{10}{12}\selectfont\setlength{\tabcolsep}{2pt}\renewcommand{\arraystretch}{1.15}\begin{tabular}{l}
Passage through perihelion 1811 Nov. 11.\\$0^{\mathrm u}\;26^{\mathrm m}\;4^{\mathrm s}$
\end{tabular}\end{center}
\begin{center}\fontsize{9}{11}\selectfont\setlength{\tabcolsep}{2pt}\renewcommand{\arraystretch}{1.15}\begin{tabular}{lcr}
Longitude of perihelion & $.\;.\;.\;.\;.\;.\;.\;.\;.\;.\;.\;.$ & $47^\circ\;27'\;27''1$\\
Longitude of the ascending node & $.\;.\;.\;.\;.\;.\;.\;.\;.$ & $93\;1\;52.1$\\
\end{tabular}\end{center}
\noindent both reckoned from the mean equinox of 1 January 1812
\begin{center}\fontsize{9}{11}\selectfont\setlength{\tabcolsep}{2pt}\renewcommand{\arraystretch}{1.15}\begin{tabular}{lcr}
Inclination of the orbit & $.\;.\;.\;.\;.\;.\;.\;.\;.\;.\;.\;.\;.\;.$ & $31^\circ\;17'\;11''0$\\
Logarithm of the least distance & $.\;.\;.\;.\;.\;.\;.$ & $0.1992359$\\
Eccentricity & $.\;.\;.\;.\;.\;.\;.\;.\;.\;.\;.\;.\;.\;.\;.\;.\;.$ & $0.98271088$\\
Semimajor axis & $.\;.\;.\;.\;.\;.\;.\;.\;.\;.\;.\;.\;.\;.$ & $91.5088$\\
Sidereal orbital period & $.\;.\;.\;.\;.\;.\;.\;.\;.\;.\;.\;.$ & $875.4\;\text{years.}$\\
\end{tabular}\end{center}

\par\indent Mr. \textsc{Nicolai} is saving a detailed account of this work for publication elsewhere. It shows the excellent agreement of this ellipse with the observations and the impossibility of reconciling them with a parabolic orbit.
